\documentclass[10pt,aspectratio=169]{beamer}
\usepackage{bma02}

\title[Convergencia absoluta e impropias]{Convergencia absoluta}
\subtitle{Integrales impropias de primera, segunda y tercera especie\\[2pt]\small Semana 13 -- Sesión 1}

\begin{document}

\begin{frame}
  \titlepage
\end{frame}

\begin{frame}{Objetivos de la sesión}
  \begin{enumerate}
    \item Distinguir convergencia \resaltar{absoluta} y \resaltar{condicional} de series; probar que la absoluta implica la convergencia.
    \item Definir las \resaltar{integrales impropias} de primera especie (intervalo infinito), segunda especie (integrando no acotado)
      y tercera especie (ambas).
    \item Calcular integrales impropias mediante límites y reconocer cuándo divergen.
    \item Aplicar a problemas de física y probabilidad.
  \end{enumerate}
  \vspace{1mm}
  {\small\textbf{Competencia 4 del sílabo:} Demuestra análisis y síntesis ejecutando cálculos en problemas aplicados a su carrera.}
  \begin{block}{Plan (2 horas; incluye verificación previa de 5 min)}
    \begin{tabular}{ll}
      0:00 -- 0:25 & Convergencia absoluta y condicional. Reordenamientos\\
      0:25 -- 1:00 & Integrales impropias de primera especie. La integral $p$\\
      1:00 -- 1:30 & Segunda especie: singularidades en extremos o interiores\\
      1:30 -- 1:45 & Tercera especie. Aplicaciones (función Gamma, velocidad de escape)\\
      1:45 -- 2:00 & Actividad colaborativa, resumen y cierre metacognitivo
    \end{tabular}
  \end{block}
\end{frame}

\begin{frame}{Aula invertida: preparación y verificación}
  \small
  \begin{previa}[antes de la sesión]
    \begin{enumerate}
      \item Lea convergencia absoluta e integrales impropias de las tres especies.
      \item Decida si converge $\int_1^\infty\frac{dx}{x^2}$ y $\int_0^1\frac{dx}{\sqrt x}$.
      \item Pregunta guía: ¿qué significa ``valor principal'' en una integral de tercera especie?
    \end{enumerate}
  \end{previa}
  \begin{block}{Verificación al inicio (5 min, en pares)}
    \begin{enumerate}
      \item ¿Para qué $p$ converge $\int_1^\infty\frac{dx}{x^p}$?
      \item ¿Qué especie es $\int_{-1}^1\frac{dx}{x}$?
    \end{enumerate}
    Los pares discuten y votan; el docente aclara solo lo que genera dudas y ajusta el ritmo de la sesión.
  \end{block}
\end{frame}

\section{Convergencia absoluta}

\begin{frame}{Convergencia absoluta y condicional}
  \begin{definicion}
    $\sum a_n$ es \resaltar{absolutamente convergente} si $\sum|a_n|$ converge; \resaltar{condicionalmente convergente}
    si $\sum a_n$ converge pero $\sum|a_n|$ diverge.
  \end{definicion}
  \begin{teorema}
    Si $\sum|a_n|$ converge, entonces $\sum a_n$ converge y $\left|\sum a_n\right|\le\sum|a_n|$.
  \end{teorema}
  \textbf{Prueba.} $0\le a_n+|a_n|\le2|a_n|$. Por comparación, $\sum(a_n+|a_n|)$ converge, y
  $\sum a_n=\sum(a_n+|a_n|)-\sum|a_n|$ es diferencia de convergentes. $\blacksquare$
  \begin{ejemplo}
    \begin{itemize}
      \item $\displaystyle\sum\frac{\sen n}{n^2}$: \ $\left|\frac{\sen n}{n^2}\right|\le\frac1{n^2}$ $\Rightarrow$ absolutamente convergente.
      \item $\displaystyle\sum\frac{(-1)^n}{\sqrt n}$: converge (Leibniz) pero $\sum\frac1{\sqrt n}$ diverge $\Rightarrow$ condicional.
      \item $\displaystyle\sum\frac{(-3)^n}{n!}$: \ $\frac{|a_{n+1}|}{|a_n|}=\frac{3}{n+1}\to0$ $\Rightarrow$ absolutamente convergente (cociente).
    \end{itemize}
  \end{ejemplo}
\end{frame}

\begin{frame}{¿Por qué importa? Reordenamientos}
  \begin{teorema}[Dirichlet / Riemann]
    \begin{itemize}
      \item Si $\sum a_n$ converge \textbf{absolutamente}, cualquier reordenamiento converge a la \emph{misma} suma.
      \item Si converge \textbf{condicionalmente}, para cualquier $L\in\R$ existe un reordenamiento que suma $L$ (¡o diverge!).
    \end{itemize}
  \end{teorema}
  \begin{ejemplo}
    $\ln2=1-\frac12+\frac13-\frac14+\frac15-\frac16+\cdots$. \ Reordenando (un positivo, dos negativos):
    \[ \left(1-\tfrac12\right)-\tfrac14+\left(\tfrac13-\tfrac16\right)-\tfrac18+\left(\tfrac15-\tfrac1{10}\right)-\tfrac1{12}+\cdots
       =\tfrac12\left(1-\tfrac12+\tfrac13-\tfrac14+\cdots\right)=\tfrac12\ln2. \]
    ¡Los mismos términos, otra suma!
  \end{ejemplo}
  \begin{idea}
    Las series absolutamente convergentes se comportan como sumas finitas; las condicionales, no.
  \end{idea}
\end{frame}

\section{Integrales impropias de primera especie}

\begin{frame}{Primera especie: intervalos no acotados}
  \begin{columns}
    \begin{column}{0.58\textwidth}
      \begin{definicion}
        Si $f$ es integrable en $[a,t]$ para todo $t>a$:
        \[ \int_a^\infty f(x)\dx=\lim_{t\to\infty}\int_a^tf(x)\dx. \]
        Si el límite existe (finito), la integral \resaltar{converge}; si no, \resaltar{diverge}. Análogo para $\int_{-\infty}^b$.
      \end{definicion}
      \begin{ejemplo}[1]
        $\displaystyle\int_1^\infty\frac{dx}{x^2}=\lim_{t\to\infty}\left[-\frac1x\right]_1^t=\lim_{t\to\infty}\left(1-\frac1t\right)=1.$
      \end{ejemplo}
      \begin{ejemplo}[2]
        $\displaystyle\int_1^\infty\frac{dx}{x}=\lim_{t\to\infty}\ln t=\infty$: \ diverge.
      \end{ejemplo}
    \end{column}
    \begin{column}{0.42\textwidth}
      \centering
      \begin{tikzpicture}
        \begin{axis}[bma, xmin=0, xmax=8.5, ymin=0, ymax=1.3, width=6.3cm, height=4.5cm, xtick={1}, ytick={1}, domain=1:8.3, samples=100]
          \addplot[fill=uniAzul!25, draw=none] {1/x^2} \closedcycle;
          \addplot[curva] {1/x^2};
          \addplot[curva2, dashed] {1/x};
          \node[uniAzul] at (axis cs:4,0.15) {\scriptsize $\frac1{x^2}$: área $1$};
          \node[uniGranate] at (axis cs:5,0.5) {\scriptsize $\frac1x$: área $\infty$};
        \end{axis}
      \end{tikzpicture}

      \small ¡Las dos curvas se parecen, pero $\frac1x$ decrece demasiado lento!
    \end{column}
  \end{columns}
\end{frame}

\begin{frame}{La integral $p$ y otros ejemplos}
  \begin{teorema}[integral $p$]
    \[ \int_1^\infty\frac{dx}{x^p}\ \text{converge}\iff p>1,\quad\text{y en ese caso vale }\frac1{p-1}. \]
  \end{teorema}
  \begin{ejemplo}[3]
    $\displaystyle\int_0^\infty e^{-ax}\dx=\lim_{t\to\infty}\frac{1-e^{-at}}{a}=\frac1a$ \ ($a>0$). \qquad
    $\displaystyle\int_0^\infty xe^{-x}\dx=\lim_{t\to\infty}\Big[-(x+1)e^{-x}\Big]_0^t=1.$
  \end{ejemplo}
  \begin{ejemplo}[4: ambos extremos infinitos]
    Se separa en un punto $c$: $\displaystyle\int_{-\infty}^\infty f=\int_{-\infty}^cf+\int_c^\infty f$ (deben converger \textbf{ambas}).
    \[ \int_{-\infty}^\infty\frac{dx}{1+x^2}=\lim_{s\to-\infty}(-\arctan s)+\lim_{t\to\infty}\arctan t=\frac\pi2+\frac\pi2=\pi. \]
  \end{ejemplo}
  \small \alert{Ojo:} $\int_{-\infty}^\infty x\dx$ diverge, aunque $\lim_{t\to\infty}\int_{-t}^tx\dx=0$ (``valor principal'', otro concepto).
\end{frame}

\begin{frame}{Aplicación: velocidad de escape}
  \begin{columns}
    \begin{column}{0.6\textwidth}
      \begin{ejemplo}[5]
        El trabajo para llevar un cohete de masa $m$ desde la superficie terrestre (radio $R$) hasta el ``infinito'' contra
        $F(r)=\dfrac{GMm}{r^2}$ es
        \[ W=\int_R^\infty\frac{GMm}{r^2}\,dr=\frac{GMm}{R}. \]
        Igualando a la energía cinética $\tfrac12mv^2$:
        \[ v_e=\sqrt{\frac{2GM}{R}}\approx11.2\ \text{km/s}. \]
      \end{ejemplo}
      \small ¡Un trabajo \emph{finito} sobre una distancia \emph{infinita}!
    \end{column}
    \begin{column}{0.4\textwidth}
      \centering
      \begin{tikzpicture}
        \begin{axis}[bma, xlabel={$r$}, ylabel={$F$}, xmin=0, xmax=6.5, ymin=0, ymax=1.3, width=6cm, height=4.6cm, xtick={1}, xticklabels={$R$}, ytick=\empty, domain=1:6.3]
          \addplot[fill=uniAzul!25, draw=none] {1/x^2} \closedcycle;
          \addplot[curva] {1/x^2};
          \node at (axis cs:3,0.5) {\scriptsize $W=\int_R^\infty F\,dr$};
        \end{axis}
      \end{tikzpicture}
    \end{column}
  \end{columns}
\end{frame}

\section{Segunda especie}

\begin{frame}{Segunda especie: integrandos no acotados}
  \begin{definicion}
    Si $f$ es continua en $(a,b]$ y no acotada cerca de $a$:\qquad
    $\displaystyle\int_a^bf(x)\dx=\lim_{s\to a^+}\int_s^bf(x)\dx$.\\
    Análogo si la singularidad está en $b$. Si está en un punto interior $c$, se separa: $\int_a^c+\int_c^b$.
  \end{definicion}
  \begin{columns}
    \begin{column}{0.55\textwidth}
      \begin{ejemplo}[6]
        $\displaystyle\int_0^1\frac{dx}{\sqrt x}=\lim_{s\to0^+}\Big[2\sqrt x\Big]_s^1=\lim_{s\to0^+}\big(2-2\sqrt s\big)=2.$
      \end{ejemplo}
      \small ¡Una región no acotada con área finita!
    \end{column}
    \begin{column}{0.45\textwidth}
      \centering
      \begin{tikzpicture}
        \begin{axis}[bma, xmin=0, xmax=1.3, ymin=0, ymax=6, width=6cm, height=4.6cm, xtick={1}, ytick={1}, samples=200]
          \addplot[fill=uniAzul!25, draw=none, domain=0.03:1] {1/sqrt(x)} \closedcycle;
          \addplot[curva, domain=0.028:1.2] {1/sqrt(x)};
          \addplot[curva2, dashed, domain=0.17:1.2] {1/x};
          \node[uniAzul] at (axis cs:0.6,2.3) {\scriptsize $\frac{1}{\sqrt x}$: área $2$};
          \node[uniGranate] at (axis cs:0.45,5) {\scriptsize $\frac1x$: área $\infty$};
        \end{axis}
      \end{tikzpicture}
    \end{column}
  \end{columns}
\end{frame}

\begin{frame}{La integral $p$ de segunda especie}
  \begin{teorema}[integral $p$, segunda especie]
    $\displaystyle\int_0^1\frac{dx}{x^p}$ converge $\iff p<1$, \ y vale $\dfrac1{1-p}$.
  \end{teorema}
  \textbf{Prueba.} Para $p\ne1$: $\displaystyle\int_s^1x^{-p}\dx=\frac{1-s^{1-p}}{1-p}$, y $s^{1-p}\to0$ si y solo si $p<1$.
  Para $p=1$: $-\ln s\to\infty$. $\blacksquare$
  \begin{center}
    \renewcommand{\arraystretch}{1.4}
    \begin{tabular}{c|c|c}
      & $\int_0^1x^{-p}\dx$ & $\int_1^\infty x^{-p}\dx$\\ \hline
      $p<1$ & converge & diverge\\
      $p=1$ & diverge & diverge\\
      $p>1$ & diverge & converge
    \end{tabular}
  \end{center}
  \begin{ejemplo}[7]
    $\displaystyle\int_0^1\ln x\dx=\lim_{s\to0^+}\Big[x\ln x-x\Big]_s^1=-1$ \ (pues $s\ln s\to0$).
  \end{ejemplo}
\end{frame}

\begin{frame}{Singularidades interiores}
  \begin{ejemplo}[8: recordando un error de la Semana 4]
    $\displaystyle\int_{-1}^1\frac{dx}{x^2}=\int_{-1}^0\frac{dx}{x^2}+\int_0^1\frac{dx}{x^2}$. \
    Pero $\displaystyle\int_0^1\frac{dx}{x^2}=\lim_{s\to0^+}\left(\frac1s-1\right)=\infty$. \ \textbf{Diverge}.
  \end{ejemplo}
  \begin{ejemplo}[9]
    $\displaystyle\int_0^3\frac{dx}{(x-1)^{2/3}}$ (singularidad en $x=1$):
    \[ \int_0^1\frac{dx}{(x-1)^{2/3}}=\lim_{t\to1^-}\Big[3(x-1)^{1/3}\Big]_0^t=3,\qquad
       \int_1^3\frac{dx}{(x-1)^{2/3}}=\lim_{s\to1^+}\Big[3(x-1)^{1/3}\Big]_s^3=3\sqrt[3]2. \]
    Converge y vale $3\big(1+\sqrt[3]2\big)$.
  \end{ejemplo}
  \begin{block}{Antes de aplicar el TFC2, verifique}
    ¿Es $f$ continua (o al menos acotada) en todo $[a,b]$? \ Si no, la integral es impropia.
  \end{block}
\end{frame}

\section{Tercera especie}

\begin{frame}{Tercera especie: ambos fenómenos}
  Intervalo no acotado \textbf{y} integrando no acotado. Se separa en integrales de una sola especie; la integral converge si
  \textbf{todas} las partes convergen.
  \begin{ejemplo}[10]
    $\displaystyle\int_0^\infty\frac{dx}{\sqrt x\,(1+x)}$. \ Con $u=\sqrt x$, $dx=2u\du$:
    \[ \int_0^\infty\frac{2u\du}{u(1+u^2)}=2\lim_{t\to\infty}\arctan t=\pi. \]
  \end{ejemplo}
  \begin{ejemplo}[11]
    $\displaystyle\int_0^\infty\frac{dx}{x^2}=\int_0^1\frac{dx}{x^2}+\int_1^\infty\frac{dx}{x^2}$: la primera diverge $\Rightarrow$ diverge.
  \end{ejemplo}
  \begin{observacion}
    $\displaystyle\int_0^\infty\frac{dx}{x^p}$ diverge para \textbf{todo} $p$: cerca de $0$ necesita $p<1$, cerca de $\infty$ necesita $p>1$.
  \end{observacion}
\end{frame}

\begin{frame}{La función Gamma}
  \begin{columns}
    \begin{column}{0.58\textwidth}
      \begin{definicion}
        Para $x>0$: \quad $\displaystyle\Gamma(x)=\int_0^\infty t^{x-1}e^{-t}\dt$.
      \end{definicion}
      Es de tercera especie si $0<x<1$ (singularidad en $t=0$). Converge para todo $x>0$ (lo probaremos con los criterios de la
      próxima sesión).
      \begin{proposicion}
        $\Gamma(1)=1$, \quad $\Gamma(x+1)=x\,\Gamma(x)$ (partes), \quad $\Gamma(n+1)=n!$.
      \end{proposicion}
      \small Generaliza el factorial: $\Gamma\left(\frac12\right)=\sqrt\pi$, así ``$\left(-\frac12\right)!=\sqrt\pi$''. Aparece en
      física estadística, mecánica cuántica y probabilidad.
    \end{column}
    \begin{column}{0.42\textwidth}
      \centering
      \begin{tikzpicture}
        \begin{axis}[bma, xmin=0, xmax=4.6, ymin=0, ymax=7, width=6cm, height=5.3cm, xtick={1,2,3,4}, ytick={1,2,6}]
          \addplot[curva, smooth] coordinates {(0.15,6.2203) (0.2,4.5908) (0.3,2.9916) (0.4,2.2182) (0.5,1.7725) (0.7,1.2981) (1,1) (1.2,0.9182) (1.46,0.8856) (1.7,0.9086) (2,1) (2.5,1.3293) (3,2) (3.5,3.3234) (4,6) (4.1,6.8126)};
          \addplot[only marks, mark=*, uniGranate] coordinates {(1,1) (2,1) (3,2) (4,6)};
          \node[uniGranate] at (axis cs:3.2,5.5) {\scriptsize $n!=\Gamma(n+1)$};
        \end{axis}
      \end{tikzpicture}
    \end{column}
  \end{columns}
\end{frame}

\begin{frame}{Actividad colaborativa: Clasificar y decidir}
  \small
  \begin{activa}[Estudio de casos]
    Cada grupo clasifica la especie y decide la convergencia de $\int_1^\infty\frac{dx}{x^p}$, $\int_0^1\frac{dx}{x^p}$ y $\int_{-1}^1\frac{dx}{x}$ (con su valor principal). Luego interpreta $\int_R^\infty\frac{GMm}{r^2}dr$ como \textbf{trabajo para escapar} del campo gravitatorio y discute por qué es finito.
  \end{activa}
  \begin{itemize}
    \item \textbf{Organización:} equipos de 3 a 4 con roles rotativos (coordinador, redactor, verificador, portavoz); 12--15 minutos.
    \item \textbf{Producto:} Tabla de especie y convergencia, y la interpretación física.
    \item \textbf{Retroalimentación:} el docente recorre los equipos, hace preguntas y no da la solución hasta la puesta en común.
  \end{itemize}
\end{frame}

\section{Ejercicios}

\begin{frame}{Ejercicios propuestos}
  \small
  \begin{ejercicio}
    \begin{enumerate}
      \item Clasifique (absoluta, condicional, divergente): $\sum\frac{(-1)^n n}{n^2+1}$, \ $\sum\frac{\cos(n\pi/3)}{n\sqrt n}$, \ $\sum(-1)^n\frac{n!}{2^n}$.
      \item Calcule o pruebe que diverge:
        \begin{multicols}{3}
        \begin{enumerate}
          \item[(a)] $\int_0^\infty\frac{x}{(1+x^2)^2}\dx$
          \item[(b)] $\int_2^\infty\frac{dx}{x\ln^2x}$
          \item[(c)] $\int_0^\infty e^{-x}\sen x\dx$
          \item[(d)] $\int_{-\infty}^\infty xe^{-x^2}\dx$
          \item[(e)] $\int_0^2\frac{dx}{\sqrt{4-x^2}}$
          \item[(f)] $\int_0^{\pi/2}\tan x\dx$
          \item[(g)] $\int_1^2\frac{dx}{x\sqrt{x^2-1}}$
          \item[(h)] $\int_{-1}^8\frac{dx}{\sqrt[3]x}$
          \item[(i)] $\int_0^\infty\frac{dx}{x^2+3x+2}$
        \end{enumerate}
        \end{multicols}
      \item Halle los valores de $p$ para los que converge $\int_e^\infty\frac{dx}{x(\ln x)^p}$.
      \item Pruebe que $\Gamma(x+1)=x\Gamma(x)$ y que $\Gamma(n+1)=n!$.
      \item \textbf{Cuerno de Gabriel}: al rotar $y=\frac1x$, $x\ge1$, alrededor del eje $x$, pruebe que el volumen es finito ($\pi$) y el área
        de la superficie $2\pi\int_1^\infty\frac1x\sqrt{1+\frac1{x^4}}\dx$ es infinita.
      \item La vida media de una partícula tiene densidad $f(t)=\lambda e^{-\lambda t}$, $t\ge0$. Pruebe que $\int_0^\infty f=1$ y halle la
        vida media $\int_0^\infty t\,f(t)\dt$.
    \end{enumerate}
  \end{ejercicio}
\end{frame}

\begin{frame}{Metacognición y uso responsable de la IA}
  \small
  \begin{columns}[t]
    \begin{column}{0.5\textwidth}
      \begin{metacog}[cierre, 3 min]
        \begin{enumerate}
          \metacomun
          \item ¿Distingo bien cada especie? ¿Qué error cometo con más frecuencia?
        \end{enumerate}
      \end{metacog}
    \end{column}
    \begin{column}{0.5\textwidth}
      \begin{iabox}
        Pida a una IA el valor principal de $\int_{-1}^1\frac{dx}{x}$. Compruebe que lo defina como límite simétrico.\\[2pt]
        \textit{Regla:} primero intente usted; luego use la herramienta para contrastar, nunca para sustituir su razonamiento. Anote el cambio en su bitácora.
      \end{iabox}
    \end{column}
  \end{columns}
\end{frame}

\begin{frame}{Resumen}
  \begin{itemize}
    \item Absoluta $\Rightarrow$ convergente. Las condicionales dependen del orden de los términos.
    \item 1.ª especie: $\int_a^\infty f=\lim_{t\to\infty}\int_a^tf$. \ $\int_1^\infty x^{-p}$ converge $\iff p>1$.
    \item 2.ª especie: $\int_a^bf=\lim_{s\to a^+}\int_s^bf$. \ $\int_0^1x^{-p}$ converge $\iff p<1$.
    \item 3.ª especie: separar; converge si todas las partes convergen.
    \item Con dos extremos problemáticos, \textbf{nunca} tomar un único límite simétrico.
  \end{itemize}
  \vspace{3mm}
  \textbf{Próxima sesión:} criterios de convergencia para integrales impropias.
\end{frame}


\begin{frame}{Evaluación y trabajo posterior}
  \small
  \begin{evalbox}[Práctica del miércoles de la semana 13: Calificada 4]
    \textbf{Cobertura:} semanas 9, 10 y 11. \textbf{Naturaleza:} práctica calificada individual escrita.
  \end{evalbox}
  \begin{itemize}
    \item \criterios
    \item \declaracionia
    \item \textbf{Portafolio:} guarde el producto de la actividad, la bitácora de metacognición y los ejercicios propuestos de esta sesión; son evidencia de progreso y base de la retroalimentación.
  \end{itemize}
\end{frame}

\end{document}
